\subsection{二重Veronese型と直線基点型の全H3排除}
\label{subsec:cubic-CY-DL-cover}
\begin{theorem}\label{thm:cubic-CY-DL-cover-empty}
零交換子CY包絡を持つH3の対称射影は、定理\ref{thm:cubic-CY-CI22-seven-projections}の
$D,L$ のどちらにもならない。
従ってこの領域に残る必要対称射影型は、有限基点の五型 $1111,211,22,31,4$ に限る。
\end{theorem}
\begin{proof}
固定出典\cite{AS4CYpinned}の二重被覆・直線基点の全係数論証から、必要正規化だけを使う。
存在十分性・全同型・自己同型群の分類はこの排除に不要であり、輸入しない。
両場合の全22混合方向と全巡回係数の完全性を元tensorで独立確認した。
以下では $R=U+KB^tA^{-1}$、$D=C_0+\gamma J-B^tA^{-1}B$ と書く。
$U,K$ は対称・交代tensor基底であり、$D=0$ がsource rank6の全Schur条件である。

\paragraph{二重Veronese型 $D$}
基底を $(a^2,ab,b^2,c^2,cd,d^2)$ とする。
rank6の印付き二次関係は、二つのVeronese関係の非零係数和である。
両二元ブロックの行列式を調整し、
$A_{02}=A_{20}=A_{35}=A_{53}=1/2$、$A_{11}=A_{44}=-1$ とできる。
全22係数 $x_0,\ldots,x_{21}$ の元混合行列に対し、Schur条件から
\[
 \gamma=0,\qquad x_3^2-4x_0x_6=-1/4,\qquad
 x_{16}^2-4x_{13}x_{21}=-1/4
\]
となる。
非退化な二元二次形式の固定判別式上の $SL_2$ 推移性により
\[
 x_0=x_6=x_{13}=x_{21}=0,\qquad x_3=x_{16}=i/2
\]
へ移せる。
元Schurの第0・第5行のrank8の線型系から
$x_4=x_5=x_7=x_8=x_{14}=x_{15}=x_{17}=x_{19}=0$ となる。
残る全式は正確に
\[
\begin{gathered}
 x_1x_9=x_2x_{10}=x_{11}x_{18}=x_{12}x_{20}=0,\\
 x_1x_{10}+x_{11}x_{20}=0,\qquad
 x_{12}x_{18}+x_2x_9=0.
\end{gathered}
\]
$(x_9,x_{10}),(x_1,x_2),(x_{18},x_{20}),(x_{11},x_{12})$ のいずれかが零なら、
実関係にそれぞれ $a^2,b^2,c^2,d^2$ が現れる。
H3包絡は整域なので、四対は全て非零である。
$x_1\ne0$ と $x_1=0$ に分ければ、残るsupportは全て
\[
\begin{array}{c|c|c}
 &\text{非零係数 }(p,q,r,z)&\text{条件}\\\hline
\mathrm I&(x_1,x_{10},x_{11},x_{20})&pq+rz=0\\
\mathrm{II}&(x_2,x_9,x_{12},x_{18})&pq+rz=0
\end{array}
\]
の二つで尽くされる。
生成元の対角変更により四係数は $(1,h,1,-h)$、$h=pq\ne0$ へ正規化できる。
必要な根は $R^2=U/p$ 又は $R^4=1/(p^3r)$ だけであり、全分母は非零である。
選んだ対角変更は二ブロックの行列式を等しくするので、正規化した $A$ を保つ。
この全正規化と元16語の関係移送を、多項式恒等式として検証した。

得られる二族は、定理\ref{thm:as4-double-cover-h3-exclusion}の二表示に文字どおり一致する。
対応する六関係基底の変換行列式も非零定数なので、全 $h\ne0$ で一致する。
同定理は全パラメータで中心三次を排除しているため、$D$ 型の全候補が消える。

\paragraph{直線基点型 $L$}
基底を $(a^2,ab,b^2,bc,ac-bd,ad)$ とし、全rank6関係の印付けを
$A_{03}=1/2,A_{14}=A_{25}=-1/2$ とその対称成分へ正規化する。
全22混合係数を残し、最後の交代列の上三成分を
\[
 v=\left(-(x_{18}+x_{21})/2,\ x_{16}-x_{20},\ (x_{13}+x_{15})/2\right)^t
\]
とすると、恒等的に $D_{15}+D_{45}=2(v_1^2-4v_0v_2)$ である。
従って $v$ の二元二次形式はrankが高々1。
$\langle a,b\rangle$ 上の $T\in GL_2$ と $\langle c,d\rangle$ 上の $T^{-t}$ の作用は
$A$ を $\det T$ 倍し、topを戻すと $v\mapsto\operatorname{Sym}^2(T)v/(\det T)^2$ となる。
この全四係数の共変恒等式を確認した。
非零平方への推移性により $v=0$ 又は $(1,0,0)$ で全場合を尽くす。

$v=0$ では $a,b$ を零にした全関係が零となり、自由二元代数への全射がAS4の増大に反する。
$v=(1,0,0)$ ではSchurの最後の列の全線型条件から
\[
 x_{12}=x_{13}=x_{15}=x_{16}=x_{20}=0,\qquad
 x_{18}=-2-x_{21},\qquad \gamma=2x_{11}
\]
となる。さらに $x_9=0$ を得る。
$z=x_{21},u=x_7,w=x_{10},p=x_{11}$ と書けば、残る必要四式は
\[
 e_1=zw=0,\quad e_2=(z+2)u=0,\quad
 e_3=p-zu+(z+2)w=0,\quad e_4=p(z+1)-2zu=0.
\]
恒等式
\[
 e_4-(z+1)e_3-(z-3)e_2+(z+3)e_1=6u-2w,\qquad
 z(6u-2w)+2e_1-6e_2=-12u
\]
から $u=w=p=0$ となる。
ここで $z$ や $z+2$ を割っておらず、全特殊値を保持している。
元関係の第0列は $a^2+2p[a,b]$ なので、実平方 $a^2$ が関係となる。
整域性に反し、$L$ 型も全排除される。
\end{proof}

公開再現器は\\
\texttt{verify\_cubic\_CY\_DL\_cover.py}。
両方の全22混合核と全cyclic条件、全D正規化と既存二族への全値関係移送、
全L共変式・全rank分岐・分母なし消去・実平方を元係数のまま検証する。
ここでの $D,L$ 全排除を、残る有限基点五型や全CY分類の完成と同一視しない。
